\chapter{Introduction}\label{ch:intro}

\section{Problem Background}

The global transition away from fossil fuels requires scalable, carbon-free energy
carriers. Molecular hydrogen (H\textsubscript{2}) is widely regarded as one of the
most attractive options because it has the highest gravimetric energy density of any
known fuel (120 MJ/kg), produces only water upon combustion, and can be stored and
transported using existing infrastructure \cite{turner2004sustainable}.
Currently, however, more than 95\% of hydrogen is produced by steam methane
reforming~— a process that emits roughly 10 kg of CO\textsubscript{2} per
kilogram of H\textsubscript{2} \cite{crabtree2004hydrogen}.

Photoelectrochemical (PEC) water splitting offers a pathway to produce green hydrogen
directly from sunlight and water using a semiconductor photoelectrode immersed in an
aqueous electrolyte \cite{fujishima1972electrochemical}.
When photons with energy exceeding the semiconductor bandgap are absorbed, they
generate electron--hole pairs.
The holes oxidise water at the photoanode (oxygen evolution reaction, OER), while
the electrons reduce protons at the cathode (hydrogen evolution reaction, HER),
as described by:

\begin{align}
  &\text{Anode (OER):}\quad 2\,\text{H}_2\text{O} \;\xrightarrow{h\nu}\; \text{O}_2 + 4\,\text{H}^+ + 4e^- \\
  &\text{Cathode (HER):}\quad 4\,\text{H}^+ + 4e^- \;\to\; 2\,\text{H}_2
\end{align}

The efficiency of this process is characterised by the Solar-to-Hydrogen (STH)
efficiency, defined as:
\begin{equation}
  \text{STH} = \frac{J_{ph} \times 1.23\,\text{V} \times \eta_F}{P_{in}} \times 100\%
  \label{eq:sth}
\end{equation}
where $J_{ph}$ is the photocurrent density (mA/cm\textsuperscript{2}),
1.23\,V is the thermodynamic minimum potential for water splitting,
$\eta_F$ is Faradaic efficiency, and
$P_{in}$ is the incident light power density (mW/cm\textsuperscript{2}).

Despite decades of research, only a handful of materials approach the theoretical
STH ceiling for their bandgap~\cite{walter2010solar}.
The search space is enormous: performance depends on the semiconductor bandgap,
nanostructure, doping, cocatalyst, electrolyte composition, pH, applied bias,
illumination conditions, and synthesis method.
A researcher must vary all of these simultaneously, making purely experimental
optimisation prohibitively slow.

\section{Motivation and Relevance}

The rise of materials informatics~\cite{ramprasad2017machine} has demonstrated
that machine learning (ML) models trained on published experimental data can
predict material properties orders of magnitude faster than density functional
theory (DFT) calculations or physical experiments.
Several recent studies have applied ML to heterogeneous catalysis
\cite{goldsmith2018machine}, battery electrolytes, and photovoltaics,
but the PEC water-splitting domain remains sparsely covered.

The main barrier is data scarcity: PEC papers typically report only a few
performance metrics, and no large, unified, open-access dataset exists.
This work addresses that gap by curating and augmenting a dataset from 529
published papers, then deploying ML models through a web application that
allows researchers to predict expected performance before committing to
synthesis.

\section{Aims and Objectives}

The aim of this project is to design, implement, and evaluate a machine
learning--based system for predicting the performance of photoelectrochemical
water-splitting devices.

\noindent The specific objectives are:

\begin{enumerate}
  \item Curate a structured dataset from published PEC literature covering
        material, structural, and operational parameters alongside measured
        performance metrics.
  \item Develop a physics-informed feature engineering pipeline that encodes
        domain knowledge (bandgap--pH interaction, overpotential, ionic
        strength, Faraday-derived H\textsubscript{2} rates) into ML-ready
        features.
  \item Implement and evaluate an ensemble of gradient-boosted and
        random-forest regressors with nested cross-validation, targeting
        three performance indicators: Photocurrent Density, STH Efficiency,
        and Hydrogen Evolution Rate.
  \item Deploy the trained models as a REST API and interactive web
        application with uncertainty-aware output (confidence score,
        model agreement).
  \item Validate that data augmentation via physics-derived estimates
        measurably improves model generalisation compared to the raw sparse
        dataset.
\end{enumerate}

\section{Structure of the Report}

Chapter~\ref{ch:A} reviews related work in PEC materials science and ML
for materials discovery.
Chapter~\ref{ch:B} describes the dataset, feature engineering, model
architecture, and system design.
Chapter~\ref{ch:C} presents experimental results and discusses model
performance.
Chapter~\ref{ch:concl} concludes with a summary of contributions and
directions for future work.
